\section{Sistemas lineales generales: Teoría general}

Dada la ecuación diferencial
\[
u' = A(t) u + B(t)
\]
donde \(u\) es una función vectorial de dimensión \(n\):
\[
    u(t) = \begin{pmatrix}
    u_1(t) \\
    \vdots \\
    u_n(t)
    \end{pmatrix}, \qquad \alpha, \beta \in \mathbb{R}, \ \alpha < \beta, \ t \in [\alpha, \beta]
\]
\(A(t)\) es una matriz de dimensión \(n \times n\) cuyas entradas son funciones continuas en el intervalo \([\alpha, \beta]\):
\[
    A(t) = \begin{pmatrix}
    a_{11}(t) & \cdots & a_{1n}(t) \\
    \vdots & \ddots & \vdots \\
    a_{n1}(t) & \cdots & a_{nn}(t)
    \end{pmatrix}, \quad a_{ij}(t) \in C([\alpha, \beta], \mathbb{K})
\]
y \(B(t)\) es una función vectorial de dimensión \(n\) (término independiente) con entradas continuas:
\[
    B(t) = \begin{pmatrix}
    b_1(t) \\
    \vdots \\
    b_n(t)
    \end{pmatrix}, \quad b_{i}(t) \in C([\alpha, \beta], \mathbb{K})
\]
al sistema $u' = A(t) u + B(t)$ se le denomina \textbf{sistema lineal general} o \textbf{sistema lineal inhomogéneo}. 

Si consideramos el caso $B(t) \equiv 0$, obtenemos el sistema:
\[
u' = A(t) u
\]
al cual se le denomina \textbf{sistema lineal homogéneo}.

Nos preguntamos cuáles son las soluciones en el intervalo \([\alpha, \beta]\) de estos sistemas. Definimos los conjuntos de soluciones como:
\[
    \Sigma_{B} = S_B = \{u \in C^{1}([\alpha, \beta], \mathbb{K}^n) : u' = A(t) u + B(t), \ t \in [\alpha, \beta]\}
\]
\[
    \Sigma_{0} = S_0 = \{u \in C^{1}([\alpha, \beta], \mathbb{K}^n) : u' = A(t) u, \ t \in [\alpha, \beta]\}
\]

\subsubsection*{Propiedades generales de los espacios de soluciones}

\begin{proposición}
    El espacio de funciones
    \[
        E \equiv C^{1}([\alpha, \beta], \mathbb{K}^n)
    \]
    es un espacio vectorial sobre el cuerpo \(\mathbb{K}\).
\end{proposición}
\begin{proof}
    Las operaciones de suma y producto por escalar se definen punto a punto. Si $u,v \in E$ y $\lambda \in \mathbb{K}$:
    \[
        (u+v)(t) = u(t) + v(t), \quad (\lambda u)(t) = \lambda u(t) \quad \forall t \in [\alpha, \beta].
    \]
    Dado que la suma de funciones derivables con continuidad es derivable con continuidad, la estructura es cerrada y cumple los axiomas de espacio vectorial.
\end{proof}

\begin{proposición}
    El conjunto de soluciones del sistema homogéneo, $\Sigma_0$, es un subespacio vectorial de $E$. Es decir, para cada $\lambda, \mu \in \mathbb{K}$ y $h_1, h_2 \in \Sigma_0$, se cumple que $\lambda h_1 + \mu h_2 \in \Sigma_0$.
\end{proposición}

\begin{proof}
    Sean $h_1, h_2 \in \Sigma_0$. Por definición, cumplen:
    \[
    h_1' = A(t) h_1, \quad h_2' = A(t) h_2, \quad \forall t \in [\alpha, \beta]
    \]
    Consideramos la combinación lineal $w = \lambda h_1 + \mu h_2$. Derivando y usando la linealidad de la matriz $A(t)$:
    \[
    w'(t) = (\lambda h_1 + \mu h_2)'(t) = \lambda h_1'(t) + \mu h_2'(t) 
    \]
    Sustituyendo las ecuaciones diferenciales:
    \[
    = \lambda A(t) h_1(t) + \mu A(t) h_2(t) = A(t) \left( \lambda h_1(t) + \mu h_2(t) \right) = A(t) w(t)
    \]
    Por lo tanto, $\lambda h_1 + \mu h_2$ es solución del sistema lineal homogéneo, lo que implica $\lambda h_1 + \mu h_2 \in \Sigma_0$.
\end{proof}

\begin{proposición}
    El conjunto de soluciones del sistema general, $\Sigma_B$, es una subvariedad afín de $E$ paralela al subespacio $\Sigma_0$.
\end{proposición}

Esta estructura se demuestra mediante las dos propiedades siguientes:

\begin{enumerate}
    \item Si tomamos dos soluciones particulares $u_1, u_2 \in \Sigma_B$, su diferencia $u_2 - u_1$ es solución del sistema homogéneo.
    \begin{proof}
        \[
            (u_2 - u_1)' = u_2' - u_1' = [A(t) u_2 + B(t)] - [A(t) u_1 + B(t)] = A(t) [u_2 - u_1].
        \]
        Luego, $u_2 - u_1 \in \Sigma_0$.
    \end{proof}
    
    \item Dado un punto $u_0 \in \Sigma_B$ (solución particular) y un vector director $h \in \Sigma_0$ (solución homogénea), la suma $u_0 + h$ es solución del sistema general.
    \begin{proof}
        \[
            (u_0 + h)' = u_0' + h' = [A(t) u_0 + B(t)] + A(t) h = A(t) [u_0 + h] + B(t).
        \]
        Luego, $u_0 + h \in \Sigma_B$.
    \end{proof}
\end{enumerate}

\subsubsection*{Existencia y Unicidad}

\begin{teorema} [Resolución del sistema lineal escalar, caso $n=1$]
    Para cada $t_0 \in [\alpha, \beta]$ y $u_0 \in \mathbb{K}$, el problema escalar
    \[
        P : \begin{cases}
            u' = a(t) u + b(t), \quad t \in [\alpha, \beta] \\
            u(t_0) = u_0
        \end{cases}
    \]
    tiene una única solución $u(t,t_0,u_0) \in C^{1}([\alpha, \beta], \mathbb{K})$. La fórmula explícita viene dada por:
    \[
        u(t) = e^{\int_{t_0}^{t} a(s) ds} u_0 + \int_{t_0}^{t} e^{\int_{s}^{t} a(r) dr} b(s) ds, \quad \forall t \in [\alpha, \beta]
    \]
\end{teorema}

\begin{proof}
    La demostración se basa en el método de \textbf{variación de las constantes}.

    Definimos la función auxiliar $h(t)$, que actúa como factor integrante, de la siguiente forma:
    \[
        h(t) := e^{\int_{t_0}^{t} a(s) ds}
    \]
    Observemos que $h(t_0) = e^0 = 1$. Además, aplicando el Teorema Fundamental del Cálculo y la regla de la cadena, su derivada es:
    \[
        h'(t) = \frac{d}{dt} \left( \int_{t_0}^{t} a(s) ds \right) e^{\int_{t_0}^{t} a(s) ds} = a(t) h(t)
    \]
    
    Proponemos una solución de la forma $u(t) = h(t)v(t)$, donde $v(t)$ es una función a determinar. Derivando respecto a $t$ obtenemos:
    \[
        u'(t) = h'(t)v(t) + h(t)v'(t)
    \]
    Sustituyendo $u$ y $u'$ en la ecuación original $u' = a(t)u + b(t)$:
    \[
        h'(t)v(t) + h(t)v'(t) = a(t) \big( h(t)v(t) \big) + b(t)
    \]
    Dado que $h'(t) = a(t)h(t)$, los términos correspondientes se simplifican, resultando en:
    \[
        h(t)v'(t) = b(t) \implies v'(t) = \frac{b(t)}{h(t)}
    \]
    Esta operación es lícita pues la función exponencial $h(t)$ no se anula en el intervalo dado.
    
    De la condición inicial $u(t_0) = u_0$, y sabiendo que $h(t_0) = 1$, deducimos que $v(t_0) = u_0$. Integrando la expresión para $v'(s)$ en el intervalo $[t_0, t]$:
    \[
        v(t) - v(t_0) = \int_{t_0}^{t} \frac{b(s)}{h(s)} ds \implies v(t) = u_0 + \int_{t_0}^{t} b(s) e^{-\int_{t_0}^{s} a(r) dr} ds
    \]
    
    Sustituyendo $v(t)$ en nuestra propuesta $u(t) = h(t)v(t)$:
    \[
        u(t) = e^{\int_{t_0}^{t} a(s) ds} \left( u_0 + \int_{t_0}^{t} b(s) e^{-\int_{t_0}^{s} a(r) dr} ds \right)
    \]
    Al distribuir el término exponencial e introducirlo en la integral, obtenemos:
    \[
        u(t) = e^{\int_{t_0}^{t} a(s) ds} u_0 + \int_{t_0}^{t} b(s) e^{\left( \int_{t_0}^{t} a(r) dr - \int_{t_0}^{s} a(r) dr \right)} ds
    \]
    Finalmente, aplicando la propiedad de aditividad de la integral ($\int_{t_0}^{t} - \int_{t_0}^{s} = \int_{s}^{t}$):
    \[
        u(t) = e^{\int_{t_0}^{t} a(s) ds} u_0 + \int_{t_0}^{t} e^{\int_{s}^{t} a(r) dr} b(s) ds
    \]
    A la ecuación anterior le llamamos \textbf{fórmula de variación de las constantes de Lagrange}.
\end{proof}

Cuando $a(t) \equiv a$ es constante, la fórmula se simplifica a la conocida:
\[
    u(t) = e^{a(t - t_0)} u_0 + \int_{t_0}^{t} e^{a(t - s)} b(s) ds
\]
En particular, si $b(t) \equiv 0$, resulta que $h(t,t_0,u_0) = e^{\int_{t_0}^{t} a(s) ds} u_0$ es la solución única del problema de Cauchy homogéneo:
\[
    \begin{cases}
    h' = a(t) h, \quad t \in [\alpha, \beta] \\
    h(t_0) = u_0
    \end{cases}
\]

\subsection{Fórmulas de variación de las constantes (de Lagrange)}

\textbf{Método de variación de las constantes}\\

\[
    u(t,t_0,u_0) = h(t,t_0,u_0) + \int_{t_0}^{t} e^{\int_{s}^{t} a(r) dr} b(s) ds, \quad \forall t \in [\alpha, \beta]
\]

\[
    \implies u(t,t_0,0) = \int_{t_0}^{t} e^{\int_{s}^{t} a(r) dr} b(s) ds = \underbrace{u(t,t_0,u_0)}_{\in \Sigma_B} + (- \underbrace{h(t,t_0,u_0)}_{\in \Sigma_0}) \in \Sigma_B 
\]
Luego la solución anterior es la solución particular del sistema inhomogéneo que anula la condición inicial. Por tanto:
\[
    u(t,t_0,u_0) = h(t,t_0,u_0) + u(t,t_0,0)
\]

\begin{teorema}
    El operador solución
    \[
        S_0 : \mathbb{K} \to \Sigma_0
    \]
    \[
        x \mapsto S_0(x) := h(t,t_0,x), \quad \forall t \in [\alpha, \beta]
    \]
    $S_0$ es la única solución de $\begin{cases}
        h' = a(t) h, \quad t \in [\alpha, \beta] \\
        h(t_0) = x
    \end{cases}$ es un isomorfismo entre espacios vectoriales. En particular $\dim(\Sigma_0) = 1$ (es una recta vectorial). Por lo tanto, el operador solución del sistema inhomogéneo
    \[
        S_B : \mathbb{K} \to \Sigma_B
    \]
    \[
        x \mapsto S_B(x) := u(t,t_0,x), \quad \forall t \in [\alpha, \beta]
    \]
    es un isomorfismo entre espacios afines. En particular, $\Sigma_B$ es una recta afín paralela a $\Sigma_0$.

\end{teorema}

\begin{proof}
    Veamos primero que $S_0$ es lineal. Sean $x,y \in \mathbb{K}$ y $\lambda, \mu \in \mathbb{K}$, entonces:
    \[
        S_0(\lambda x + \mu y) = e^{\int_{t_0}^{t} a(s) ds} (\lambda x + \mu y) = \lambda e^{\int_{t_0}^{t} a(s) ds} x + \mu e^{\int_{t_0}^{t} a(s) ds} y = \lambda S_0(x) + \mu S_0(y)
    \]
    Además, $S_0$ es inyectiva, pues si $S_0(x) = S_0(y)$, por ser $S_0$ lineal, se cumple que $S_0(x - y) = 0$, lo que implica que $x - y = 0$, es decir, $x = y$.
    \[
        S_0(z) = 0 \implies e^{\int_{t_0}^{t} a(s) ds} z = 0 \implies z = 0
    \]
    Finalmente, $S_0$ es sobreyectiva, ya que para cualquier $h \in \Sigma_0$, tenemos que $h'(t) = a(t) h(t), \ \forall t \in [\alpha, \beta]$, $h(t_0) = x$ para algún $x \in \mathbb{K}$. Por lo tanto, $h = S_0(x)$.
    \[
        h(t) = e^{\int_{t_0}^{t} a(s) ds} x = S_0(x)
    \]
\end{proof}

Las soluciones de $\Sigma_0$ son todos los puntos de la recta vectorial generada por el vector $e^{\int_{t_0}^{t} a(s) ds}$, es decir:
\[
    \Sigma_0 = \{ \lambda e^{\int_{t_0}^{t} a(s) ds} : \lambda \in \mathbb{K} \}
\]
Por otro lado, las soluciones de $\Sigma_B$ son todos los puntos de la recta afín paralela a $\Sigma_0$ que pasa por el punto $u(t,t_0,0)$, es decir:
\[
    \Sigma_B = \{ u(t,t_0,0) + \lambda e^{\int_{t_0}^{t} a(s) ds} : \lambda \in \mathbb{K} \}
\]

\ejemplo{
    Sea el problema de valor inicial definido por la ecuación diferencial lineal de primer orden $u' = au + b$, con $a, b \in \mathbb{R}$, $a \neq 0$ y la condición inicial $u(t_0) = u_0$. Para resolverlo, consideramos primero la ecuación diferencial homogénea asociada $u' - au = 0$, cuya solución se obtiene mediante la integración directa, resultando en $u_h(t) = \lambda e^{at}$, donde $\lambda$ es una constante real. Posteriormente, buscamos una solución particular $u_p(t)$ para la ecuación completa. Dado que los coeficientes son constantes, proponemos una solución de equilibrio de la forma $u_p(t) = k$. Al sustituir en la ecuación original, obtenemos $0 = ak + b$, lo que implica que la solución particular es $u_p = -b/a$. 

    La solución general de la ecuación diferencial se construye como la suma de la solución homogénea y la solución particular, resultando en la expresión $u(t) = \lambda e^{at} - \frac{b}{a}$. Para determinar el valor de la constante $\lambda$ que satisface la condición inicial específica $u(t_0) = u_0$, evaluamos la función en $t = t_0$:
    \[
    u_0 = \lambda e^{at_0} - \frac{b}{a}
    \]
    Al despejar el término que contiene la constante, obtenemos $\lambda e^{at_0} = u_0 + \frac{b}{a}$, y multiplicando ambos miembros por $e^{-at_0}$ hallamos que $\lambda = e^{-at_0} \left( u_0 + \frac{b}{a} \right)$. Finalmente, sustituimos este valor de $\lambda$ en la solución general:
    \[
    u(t) = e^{-at_0} \left( u_0 + \frac{b}{a} \right) e^{at} - \frac{b}{a}
    \]
    Aplicando las propiedades de los exponentes para combinar los términos con base $e$, llegamos a la expresión final del problema de valor inicial:
    \[
    u(t) = e^{a(t - t_0)} \left( u_0 + \frac{b}{a} \right) - \frac{b}{a}
    \]
}

\ejemplo{
    \[
        u(t) = au + e^{wt}, \quad w \in \mathbb{K}, \quad a \neq w
    \]
    \[
        h' = ah \implies h(t) = e^{at}, \quad \forall t \in \mathbb{R}
    \]
    Solución general:
    \[
        u(t) = e^{at} \lambda + p(t), \quad \lambda \in \mathbb{K}
    \]
    \[
        p(t) = me^{wt} \implies p'(t) = wme^{wt} = ame^{wt} + e^{wt} \implies m(w - a) = 1 \implies m = \frac{1}{w - a}
    \]
    Luego
    \[
    p(t) = \frac{e^{wt}}{w - a}
    \]
    resuelve la ecuación inhomogénea. Por tanto, la solución general es:
    \[
        u(t) = e^{at} \lambda + \frac{e^{wt}}{w - a}, \quad \lambda \in \mathbb{K}
    \]
    \[
        e^{at_0} = u_0 - \frac{e^{wt_0}}{w - a} \implies \lambda = e^{-at_0} \left( u_0 - \frac{e^{wt_0}}{w - a} \right)
    \]
    \[
        u(t) = e^{a(t - t_0)} \left( u_0 - \frac{e^{wt_0}}{w - a} \right) + \frac{e^{wt}}{w - a}
    \]
}

\ejemplo{
    \[
        u' = au + e^{at}, \quad a \in \mathbb{K}
    \]
    La ecuación homogénea asociada es $h' = ah$, cuya solución general es:
    \[
        h(t) = e^{at}, \quad \forall t \in \mathbb{R}
    \]
    Por tanto, la solución general del sistema inhomogéneo es:
    \[
        u(t) = e^{at} \lambda + p(t), \quad \lambda \in \mathbb{K}
    \]
    \[
        p(t) = me^{at} 
    \]
    como es solución de la homogénea, no puede ser solución del inhomogéneo. Hacemos el cambio de variable:
    \[
        u(t) = e^{at} v(t) \implies u'(t) = a e^{at} v(t) + e^{at} v'(t)
    \]
    \[
        ae^{at} v(t) + e^{at} v'(t) = a e^{at} v(t) + e^{at} \implies v'(t) = 1 \implies v(t) = t + \lambda, \quad \lambda \in \mathbb{K}
    \]
    Luego
    \[
        u(t) = e^{at} (t + \lambda) = \underbrace{\lambda e^{at}}_{\in \Sigma_0} + \underbrace{t e^{at}}_{\in \Sigma_B}, \quad \lambda \in \mathbb{K} 
    \]
    \[
        p(t) = t e^{at} \text{ es solución particular del inhomogéneo}
    \]
    \[
        u_0 = u(t_0) = \lambda e^{at_0} + t_0 e^{at_0} \implies \lambda = e^{-at_0} (u_0 - t_0 e^{at_0})
    \]
    \[
        u(t,t_0,u_0) = e^{a(t - t_0)} (u_0 - t_0 e^{at_0}) + t e^{at}
    \]
}


\begin{teorema} [Teorema Fundamental]
    Sean $\alpha < \beta$ reales, y $a_{ij},b_j \in C([\alpha, \beta], \mathbb{K})$, para $i,j = 1, \ldots, n$. Sean $A(t) = (a_{ij}(t))_{i,j=1}^{n}$ y $B(t) = (b_j(t))_{j=1}^{n}$. Entonces, para cada $t_0 \in [\alpha, \beta]$ y cada $u_0 \in \mathbb{K}^n$, el problema de valores iniciales de Cauchy:
    \[
        P : \begin{cases}
            u' = A(t) u + B(t), \quad t \in [\alpha, \beta] \\
            u(t_0) = u_0
        \end{cases}
    \]
    tiene una única solución $u(t,t_0,u_0) \in C^{1}([\alpha, \beta], \mathbb{K}^n)$.    
\end{teorema}

\begin{proof}
    Supongamos que $u(t)$ resuelve el problema $P$
    \[
        u'(s) = A(s) u(s) + B(s), \quad s \in [\alpha, \beta]
    \]
    Integrando entre $t_0$ y $t$, coordenada a coordenada, obtenemos:
    \[
        \int_{t_0}^{t} u'(s) ds = \int_{t_0}^{t} A(s) u(s) ds + \int_{t_0}^{t} B(s) ds
    \]
    \[
        u(t) - u(t_0) = \int_{t_0}^{t} A(s) u(s) ds + \int_{t_0}^{t} B(s) ds
    \]
    \[
        u(t) = u_0 + \int_{t_0}^{t} (A(s) u(s) + B(s)) ds
    \]
\end{proof}

\begin{lema}
    Supongamos que $u \in C([\alpha, \beta], \mathbb{K}^n)$. Entonces las siguientes afirmaciones son equivalentes:
    \begin{enumerate}
        \item $u \in C^1([\alpha, \beta], \mathbb{K}^n)$ resuelve $P$ en $[\alpha, \beta]$.
        \item $u(t) = u_0 + \int_{t_0}^{t} (A(s) u(s) + B(s)) ds$, para todo $t \in [\alpha, \beta]$.
    \end{enumerate}
\end{lema}

\begin{proof}
    \leavevmode
    \begin{enumerate}
        \item[$(1) \implies (2)$] Ya se ha demostrado en el teorema anterior.
        \item[$(2) \implies (1)$] Por el Teorema Fundamental del Cálculo Integral, sabemos que si $a : [\alpha, \beta] \to \mathbb{R}$ es continua
        \[
            \varphi (t) = \int_{t_0}^{t} a(s) ds \implies \varphi' (t) = a(t), \quad \varphi \in C^1([\alpha, \beta], \mathbb{R}), \quad \forall t \in [\alpha, \beta]
        \]
        \[
            \varphi'(t) = \lim_{h \to 0} \frac{\varphi(t+h) - \varphi(t)}{h} = \lim_{h \to 0} \frac{1}{h} \int_{t}^{t+h} a(s) ds 
        \]
        \[
            h \min_{s \in [t, t+h]} a(s) \leq \int_{t}^{t+h} a(s) ds \leq h \max_{s \in [t, t+h]} a(s)
        \]
        \[
            a(t) = \lim_{h \to 0} \min_{s \in [t, t+h]} a(s) \leq \lim_{h \to 0} \frac{1}{h} \int_{t}^{t+h} a(s) ds \leq \lim_{h \to 0} \max_{s \in [t, t+h]} a(s) = a(t)
        \]
        \[
            \implies \varphi'(t) = a(t)
        \]
        En el caso complejo, $a : [\alpha, \beta] \to \mathbb{C}$ continua, se descompone en sus partes real e imaginaria $a = a_1 + i a_2$, con $a_1, a_2 : [\alpha, \beta] \to \mathbb{R}$ continuas. Entonces:
        \[
            \varphi(t) = \int_{t_0}^{t} a(s) ds = \int_{t_0}^{t} a_1(s) ds + i \int_{t_0}^{t} a_2(s) ds
        \]
        \[
            \varphi'(t) = \frac{d}{dt} \left( \int_{t_0}^{t} a_1(s) ds + i \int_{t_0}^{t} a_2(s) ds \right) = \frac{d}{dt} \left( \int_{t_0}^{t} a_1(s) ds \right) + i \frac{d}{dt} \left( \int_{t_0}^{t} a_2(s) ds \right)
        \] 
        \[
            = a_1(t) + i a_2(t) = a(t)
        \]
        \[
            a : [\alpha, \beta] \to \mathbb{K}^n
        \]
        \[
            a = \begin{pmatrix}
            a_1 \\
            \vdots \\
            a_n
            \end{pmatrix}, \quad a_j : [\alpha, \beta] \to \mathbb{K} \text{ continua}, \quad j = 1, \ldots, n
        \]
        \[
            \frac{d}{dt} \left( \int_{t_0}^{t} \begin{pmatrix}
            a_1(s) \\
            \vdots \\
            a_n(s)
            \end{pmatrix} ds \right) = \begin{pmatrix}
            \frac{d}{dt} \left( \int_{t_0}^{t} a_1(s) ds \right) \\
            \vdots \\
            \frac{d}{dt} \left( \int_{t_0}^{t} a_n(s) ds \right)
            \end{pmatrix} = \begin{pmatrix}
            a_1(t) \\
            \vdots \\
            a_n(t)
            \end{pmatrix} = a(t)
        \]
    \end{enumerate}
\end{proof}